\documentclass[11pt]{article}
\usepackage{fontspec,amsmath,amssymb,amsthm,hyperref}
\usepackage[margin=1in]{geometry}
\setlength{\emergencystretch}{3em}
\title{The actual KLMC2 chain: a bound for every initial transport plan}
\author{LangevinLean mathematical exposition}
\date{}
\begin{document}\maketitle
\noindent Source module: \texttt{KineticLangevin/CanonicalKLMC2Bound.lean}. Compile this proof document with XeLaTeX or LuaLaTeX. The accompanying HTML pairs the mathematical argument with the exact Lean statements and records the scope of the formalization.


The reference is the constructed kinetic diffusion started from the normalized Gibbs target, with an independent Gaussian velocity and canonical Brownian source. The numerical chain is the implemented four-kernel KLMC2 update on the same source. All local equations, Gaussian marginals and stationary reference laws have been proved before the scalar recurrence is applied.


\section{Keep the accumulated contraction factor nonnegative}


The source restrictions imply gamma h at most one fifth and m at most gamma squared. Hence mh/(4 gamma) is at most one twentieth, so q=1-mh/(4 gamma) is nonnegative. This elementary check is needed when comparing the actual initial error and projecting the iterated recurrence.



\noindent Lean declarations: \texttt{secondOrder\_contraction\_factor\_nonneg}.

\section{Iterate the actual canonical coupling error}


For every finite-cost initial position coupling, construct the actual numerical and reference chains with shared initial Gaussian velocity and shared canonical Brownian intervals. Let e(k) be the actual transformed L2 phase discrepancy and d(k) the actual L2 position discrepancy. The native canonical one-step theorem supplies the repaired local recurrence, including the fresh stochastic Hessian terms. The cutoff condition M2 sqrt(alpha)h at most m/4 absorbs its position-dependent correction into q. The checked affine recurrence and phase projection then give d(n) at most sqrt(2)q\textasciicircum{}n times the initial position cost, plus sqrt(2)h\textasciicircum{}2 C/m and 8(M-m)h exp(-(alpha-p)/8)/m. Here C=M2 sqrt(p\textasciicircum{}2+2p)+M sqrt(Mp). Equal initial velocities identify the initial transformed norm exactly, and the constructed initial-extension variables identify the initial L2 position norm with the supplied transport plan cost. There is no stationary-law, local-error, Brownian-independence or numerical-law premise in this final per-coupling theorem. Taking the infimum over actual initial couplings is handled by the subsequent Wasserstein endpoint.



\noindent Lean declarations: \texttt{klmc2\_coupling\_position\_bound}.

\end{document}
